\documentclass{article}

\usepackage[utf8]{inputenc} 
\usepackage{amsmath}        
\usepackage{array}
\usepackage{booktabs}
\usepackage{multirow}
\usepackage{geometry}
\usepackage{caption}
\geometry{a4paper, margin=1in}


\begin{document}

\begin{table}[htbp]
  \centering
  \caption{Matrice des utilités totales $u_p(e) + u_e(p)$}
  \label{tab:matrice_utilites}
  
  \begin{tabular}{c | cccccccccc}
    \toprule
    \multirow{2}{*}{Etudiants} & \multicolumn{10}{c}{Parcours} \\
    \cmidrule(lr){2-11}
       & 0 & 1 & 2 & 3 & 4 & 5 & 6 & 7 & 8 & 9 \\
    \midrule
    0  & 6 & 5 & 6 & 9 & 10 & 20 & 18 & 10 & 16 & 12 \\
    1  & 12 & 6 & 8 & 7 & 7 & 20 & 19 & 9 & 14 & 8 \\
    2  & 8 & 2 & 6 & 3 & 9 & 5 & 4 & 7 & 9 & 5 \\
    3  & 12 & 8 & 12 & 9 & 14 & 17 & 17 & 11 & 14 & 16 \\
    4  & 12 & 17 & 10 & 8 & 9 & 14 & 15 & 13 & 10 & 8 \\
    5  & 18 & 12 & 13 & 12 & 12 & 8 & 7 & 15 & 12 & 17 \\
    6  & 5 & 3 & 11 & 8 & 5 & 16 & 12 & 19 & 12 & 3 \\
    7  & 20 & 13 & 12 & 15 & 11 & 6 & 5 & 21 & 11 & 14 \\
    8  & 4 & 4 & 7 & 5 & 4 & 9 & 9 & 9 & 6 & 8 \\
    9  & 11 & 12 & 19 & 14 & 4 & 10 & 8 & 10 & 9 & 18 \\
    10 & 10 & 7 & 3 & 8 & 20 & 5 & 10 & 2 & 8 & 5 \\
    11 & 10 & 14 & 16 & 16 & 16 & 7 & 13 & 12 & 1 & 17 \\
    12 & 15 & 13 & 12 & 6 & 9 & 3 & 15 & 16 & 8 & 14 \\
    \bottomrule
  \end{tabular}
\end{table}

\noindent\textbf{Modélisation de la recherche de l'affectation maximisant l'utilité totale (efficacité)}

\vspace{0.3cm}
\noindent\textbf{Variables:} Pour tout étudiant $i \in \{1,\dots,n\}$ et tout parcours $j \in \{1,\dots,m\}$, on définit la variable $x_{ij}$ telle que $x_{ij} = 1$ si l'étudiant $e_i$ est affecté au parcours $p_j$ et $x_{ij} = 0$ sinon.

\vspace{0.3cm}
\noindent\textbf{Fonction objectif :} 
\[ \max \sum_{i=1}^{n} \sum_{j=1}^{m} x_{ij}\left(u_{e_i}(p_j) + u_{p_j}(e_i)\right) \]

\noindent\textbf{Contraintes:}
\begin{align*}
& \sum_{j=1}^{m} x_{ij} = 1, && \forall i \in \{1,\dots,n\} && \text{(chaque étudiant est affecté à un seul parcours)} \\
& \sum_{i=1}^{n} x_{ij} \leq c_j, && \forall j \in \{1,\dots,m\} && \text{(la capacité $c_j$ de chaque parcours est respectée)} \\
& x_{ij} \in \{0, 1\}, && \forall i, j && \text{(domaine des variables)}
\end{align*}

\end{document}